\chapter*{Glosario}
\addcontentsline{toc}{chapter}{Glosario}

{\sloppy
\begin{description}

  \item[Abundancia de rocas] Proporción de la superficie ocupada por rocas por
    encima de un cierto tamaño, expresada habitualmente como porcentaje del área
    total. En estudios de selección de sitio de aterrizaje, una abundancia baja
    indica terreno relativamente limpio y valores altos señalan zonas peligrosas
    para un módulo o un rover.

  \item[Anotación colaborativa] Anotación realizada por voluntarios no
    especialistas. En AI4Mars, cada voluntario trazó polígonos sobre la imagen para
    marcar el tipo de terreno, y los polígonos de varios voluntarios se fusionaron
    en una máscara por píxel exigiendo consenso entre ellos.

  \item[Apertura morfológica] Erosión seguida de dilatación con el mismo elemento
    estructurante. Elimina manchas pequeñas y ruido aislado sin alterar
    sensiblemente las regiones grandes.

  \item[Área de una componente] Número de píxeles que pertenecen a una misma
    región conectada. Es una medida directa del tamaño proyectado del objeto.

  \item[Arena] Clase de la escala de navegación (\textit{sand}) que agrupa los
    depósitos de arena suelta.

  \item[Bandera de calidad] Etiqueta que resume la aptitud de cada escena para los
    indicadores: \texttt{ok}, \texttt{no\_bigrock}, \texttt{no\_rock},
    \texttt{mostly\_null} y \texttt{empty} (Tabla~\ref{tab:banderas}).

  \item[\textit{Bootstrap} por sol] Remuestreo que toma soles completos, con todas
    sus escenas, en lugar de escenas sueltas, para que los intervalos tengan en
    cuenta que las imágenes de un mismo sol no son independientes.

  \item[Caja envolvente] Rectángulo alineado con los ejes de la imagen que contiene
    por completo una componente conectada. De ella se derivan el ancho, el alto y
    la relación de aspecto.

  \item[Cierre morfológico] Dilatación seguida de erosión. Rellena huecos pequeños
    y une regiones separadas por una distancia mínima.

  \item[Cobertura de roca visible] Porcentaje de píxeles etiquetados de una imagen
    que corresponden a lecho rocoso o roca grande. Resume cuán rocoso es el terreno
    de la escena tal como lo registra la anotación.

  \item[Componente conectada] Conjunto de píxeles activos de una imagen binaria en
    el que cada píxel se alcanza desde los demás mediante pasos entre vecinos.
    Con conectividad de ocho vecinos se admiten también los pasos diagonales.

  \item[Correlación de distancia] Medida de dependencia entre dos variables que vale
    cero si y solo si son independientes, de modo que detecta también relaciones no
    lineales y no monótonas.

  \item[Desplazamiento circular, prueba por] Contraste de la relación entre dos
    variables que ordena las escenas en el tiempo y desplaza una serie respecto de
    la otra: cada serie conserva su autocorrelación y solo se rompe la relación
    entre ambas.

  \item[Diagrama de Bland--Altman] Gráfico de la diferencia entre dos mediciones
    frente a su media, con el error medio y los límites de acuerdo. Muestra si una
    medición se desvía de la otra de forma sistemática y cuánto.

  \item[Distribución tamaño--frecuencia] Relación entre el tamaño de las rocas y su
    número. En superficies planetarias suele ser decreciente: predominan las rocas
    pequeñas y la frecuencia cae al aumentar el tamaño.

  \item[División de aguas] Técnica de segmentación que interpreta una imagen —o una
    transformación de ella— como un relieve topográfico que se inunda desde sus
    mínimos; las regiones crecen hasta encontrarse y definen las fronteras. Aquí se
    inunda el negativo de la transformada de distancia, partiendo de semillas
    situadas en sus máximos, para separar rocas que aparecen pegadas.

  \item[Elemento estructurante] Máscara pequeña que define la vecindad empleada por
    una operación morfológica.

  \item[Fracción etiquetada] Proporción de los píxeles de una escena que tienen
    clase asignada en la máscara. Es el denominador de la cobertura sobre píxeles
    etiquetados y un indicador de calidad de cada escena.

  \item[Información mutua] Cantidad de información, en bits, que una variable aporta
    sobre otra. Vale cero si son independientes y no presupone ninguna forma de
    relación.

  \item[Instrumento común] Medición aplicada por igual a dos muestras para separar lo
    que se debe a las imágenes de lo que se debe a la anotación. En este trabajo es
    el segmentador entrenado, aplicado a la región anotable de cada escena.

  \item[Kappa ponderado] Coeficiente de acuerdo entre dos clasificaciones ordinales,
    corregido por el acuerdo esperable por azar, que penaliza más los desacuerdos
    entre categorías alejadas.

  \item[Lecho rocoso] Roca expuesta en continuidad con el sustrato, frente a los
    bloques sueltos apoyados sobre el regolito. En la escala de navegación
    corresponde a la clase \textit{bedrock}. En la práctica, en AI4Mars la clase
    incluye también muchos bloques sueltos (Sección~\ref{sec:diagnostico}).

  \item[Máscara binaria] Imagen de dos valores en la que los píxeles activos
    señalan la clase de interés y el resto el fondo.

  \item[Máscara colaborativa] Máscara producida por anotación colaborativa. En
    AI4Mars son las máscaras del conjunto de entrenamiento, por lo que el texto las
    llama también máscaras de entrenamiento.

  \item[Máscara de especialista] Máscara elaborada por especialistas de la misión
    para evaluar modelos; el texto la llama también máscara de experto. AI4Mars la
    publica en tres niveles, según cuántos especialistas deben coincidir en cada
    píxel.

  \item[Máscara de segmentación por píxel] Imagen del mismo tamaño que la original
    en la que cada píxel guarda un entero que identifica la clase de terreno a la
    que pertenece.

  \item[Máximo de prominencia] Máximo local cuya altura sobre el entorno supera un
    umbral dado. Filtrar por prominencia descarta las crestas de poca altura que
    subdividen una región continua.

  \item[Píxel sin etiqueta] Píxel que la máscara deja sin clase asignada, por falta
    de acuerdo entre anotadores o por corresponder al cuerpo del rover, al cielo o
    a distancias mayores de treinta metros. Se excluye del denominador de la
    cobertura sobre píxeles etiquetados, pero cuenta en la cobertura sobre la
    imagen completa y en la fracción etiquetada.

  \item[Población de E1 y de E2] Escenas sobre las que se calcula cada indicador:
    la cobertura, sobre todas las que tienen algún píxel etiquetado; el conteo,
    sobre las de bandera \texttt{ok}, con roca grande y sin estar casi vacías.

  \item[Reglas heurísticas de priorización] Reglas con umbrales explícitos que
    ordenan las escenas para su revisión a partir de los indicadores. No están
    calibradas para decisiones de navegación (Anexo~\ref{app:extension}).

  \item[Relación de aspecto] Cociente entre el lado mayor y el lado menor de la
    caja envolvente. Sirve para descartar formas excesivamente alargadas, poco
    plausibles como roca aislada.

  \item[Reloj de nave] Contador de segundos de a bordo que el identificador de cada
    imagen incluye. Permite ordenar las imágenes en el tiempo.

  \item[Reproducibilidad] Propiedad de un análisis que permite obtener los mismos
    resultados a partir de los mismos datos siguiendo el procedimiento descrito,
    con el código, los parámetros y las versiones de biblioteca documentados.

  \item[Roca grande] Clase de la escala de navegación (\textit{big rock}) destinada
    a los bloques discretos apoyados sobre el terreno. Es la única que el conteo
    considera.

  \item[Secuencia de adquisición] Orden en que el rover tomó las imágenes, según el
    reloj de nave. Es un orden temporal: no equivale a la distancia recorrida.

  \item[Segmentación semántica] Tarea de asignar a cada píxel de una imagen una
    etiqueta de clase, sin distinguir entre instancias de la misma clase.

  \item[Semilla o marcador] Punto o región de partida a partir de la cual crece un
    dominio en la división de aguas. Cada semilla origina, en principio, una región
    de salida.

  \item[Sobresegmentación] División de un mismo objeto en varias regiones. En el
    conteo, se produce cuando la división de aguas parte en varios trozos una
    región que corresponde a una sola roca o a un afloramiento continuo.

  \item[Sol] Día solar marciano, de unas 24 horas y 40 minutos.

  \item[Solidez] Cociente entre el área de una región y el área de su envolvente
    convexa. Valores próximos a uno indican formas convexas; valores bajos señalan
    contornos cóncavos.

  \item[Subconteo] Conteo inferior al número de rocas que un observador distingue.

  \item[Suelo] Clase de la escala de navegación (\textit{soil}) que agrupa el
    regolito de grano fino que no es arena.

  \item[Transformada de distancia] Imagen en la que cada píxel activo guarda la
    distancia al píxel de fondo más cercano. Sus máximos locales se sitúan en los
    centros de las regiones y sirven como semillas naturales.

\end{description}
}
\clearpage
